\documentclass[aspectratio=169]{beamer}
\usepackage{fontspec}

% Tipografías de la Identidad Visual
\setmainfont{Inter}
\setsansfont{Inter}
\setmonofont{JetBrains Mono}
\newfontfamily\titlefont{Montserrat}[BoldFont={Montserrat Bold}]

% Colores de la Identidad Visual
\definecolor{AzulMarino}{HTML}{1A365D}
\definecolor{GrisPizarra}{HTML}{4A5568}
\definecolor{BlancoNieve}{HTML}{F7FAFC}
\definecolor{VerdeAzulado}{HTML}{319795}

% Aplicación de colores a Beamer
\setbeamercolor{palette primary}{bg=AzulMarino, fg=BlancoNieve}
\setbeamercolor{palette secondary}{bg=VerdeAzulado, fg=BlancoNieve}
\setbeamercolor{palette tertiary}{bg=GrisPizarra, fg=BlancoNieve}
\setbeamercolor{structure}{fg=AzulMarino}
\setbeamercolor{background canvas}{bg=BlancoNieve}
\setbeamercolor{normal text}{fg=GrisPizarra}

\setbeamerfont{title}{family=\titlefont, series=\bfseries, size=\Huge}
\setbeamerfont{frametitle}{family=\titlefont, series=\bfseries, size=\LARGE}

\usepackage[utf8]{inputenc}
\usepackage{graphicx}
\usepackage{listings}

\lstset{
  basicstyle=\ttfamily\color{GrisPizarra}\small,
  keywordstyle=\color{VerdeAzulado}\bfseries,
  stringstyle=\color{AzulMarino},
  commentstyle=\color{GrisPizarra}\itshape,
  backgroundcolor=\color{BlancoNieve},
  frame=single,
  rulecolor=\color{GrisPizarra}
}

\title{Práctica Guiada: Pandas y Python}
\subtitle{Machine Learning (Laboratorio)}
\author{Clase Práctica - Semana 1}
\date{}

\begin{document}

\begin{frame}
    \titlepage
\end{frame}

\begin{frame}[fragile]{Importando Librerías y Cargando Datos}
    En Ciberseguridad, antes de bloquear un paquete (Malware), debemos leer los registros de red.
    
\begin{lstlisting}[language=Python]
import pandas as pd
from sklearn.datasets import load_iris

# Carga del dataset simulado
iris_raw = load_iris()
df = pd.DataFrame(data=iris_raw.data, 
                  columns=iris_raw.feature_names)
\end{lstlisting}
    
    \begin{itemize}
        \item \texttt{pandas}: Librería clave para analizar tablas de datos bidimensionales.
        \item \texttt{load\_iris}: Carga nuestro dataset clásico para exploración.
    \end{itemize}
\end{frame}

\begin{frame}[fragile]{Filtros y Consultas (Queries)}
    ¿Cómo filtramos los ataques maliciosos de un conjunto de 10 millones de filas?
    
\begin{lstlisting}[language=Python]
# Consulta usando query() de Pandas
filtro = df.query("species == 0 and `petal width (cm)` <= 0.2")
print(len(filtro))
\end{lstlisting}
    
    \begin{itemize}
        \item \textbf{\textcolor{VerdeAzulado}{Interpretación de Data:}} Usamos backticks (\`) si las columnas tienen espacios en su nombre.
        \item Los operadores lógicos como \texttt{and} nos permiten cruzar atributos para detectar anomalías.
    \end{itemize}
\end{frame}

\begin{frame}{Toma de Decisiones: El Negocio Detrás del Filtro}
    \begin{alertblock}{Métricas de Ciberseguridad}
    Si creamos un filtro muy estricto:
    \begin{itemize}
        \item Incrementamos los \textbf{Falsos Positivos}: Bloqueamos a los diseñadores gráficos.
        \item Pero reducimos los \textbf{Falsos Negativos}: Disminuye el riesgo de entrada de Malware.
    \end{itemize}
    \end{alertblock}
    
    ¡A programar en Jupyter!
\end{frame}

\end{document}
